%%%%%%%%%%%%%%%%%%%%%%%%%%%%%%%%%
%
%       EXERCÍCIO
%
%%%%%%%%%%%%%%%%%%%%%%%%%%%%%%%%%

\definevalorquestao

\begin{exercicioBanco}[\valorquestao]
Um posto de saúde registrou a quantidade de vacinas aplicadas contra febre amarela nos últimos cinco meses:
\begin{itemize}
    \begin{multicols}{3}
        \item 1º mês: 21;
        \item 2º mês: 22;
        \item 3º mês: 25;
        \item 4º mês: 31;
        \item 5º mês: 21.
    \end{multicols}
\end{itemize}

No início do primeiro mês, esse posto de saúde tinha 228 vacinas em estoque. A política de reposição prevê a aquisição de novas vacinas, no início do sexto mês, de tal forma que a quantidade inicial em estoque para os próximos meses seja igual a 12 vezes a média das quantidades mensais aplicadas nos últimos cinco meses.

Qual a quantidade de vacinas que deve ser adquirida no início do sexto mês?
\end{exercicioBanco}

